% ============================================================
% 02 -- SERVICE OVERVIEW
% ============================================================

\section{Service Overview: HTTP, but Make It TeX}
\label{sec:service-overview}

Before looking at the individual vulnerabilities, it helps to understand
how a request moves through LAMP.

The central application entry point was \code{/srv/main.tex}. It configured
the service's Redis and MariaDB connections, loaded the custom HTTP
implementation, installed the authentication middleware, and finally started
request processing with:

\begin{lstlisting}
\httpMiddlewareUse{\authMiddleware}
\httpServe{}
\end{lstlisting}

In other words, XeLaTeX was not merely generating a page after some other
application had handled the request.

It was part of the application itself.


% ------------------------------------------------------------
% 2.1
% ------------------------------------------------------------

\subsection{Request Flow}

At a high level, a request followed this path:

\begin{enumerate}

  \item an HTTP request reached the LAMP service;

  \item the request was passed into the XeLaTeX process;

  \item the custom HTTP code parsed the request line, headers, cookies,
        query parameters, and POST body;

  \item the requested page was selected and loaded as TeX;

  \item application state was read from Redis, MariaDB, or files under
        \code{/storage};

  \item XeLaTeX expanded the resulting macros and generated the HTTP
        response.

\end{enumerate}

The important part is that HTTP state eventually became TeX state.

Headers could become control sequences. POST parameters became macros.
Database values were read back into TeX. Files containing user-controlled
data were later included by the renderer.

This created several opportunities for data that was originally supposed
to be passive to become executable TeX.


% ------------------------------------------------------------
% 2.2
% ------------------------------------------------------------

\subsection{Where LAMP Stored Its State}

LAMP split its state across three main locations.

\begin{center}
\begin{tabular}{p{0.22\textwidth} p{0.66\textwidth}}
  \toprule
  \textbf{Storage} & \textbf{Purpose} \\
  \midrule

  Redis &
  Sessions, usernames, passwords, and temporary authentication state. \\[0.4em]

  MariaDB &
  Ships and their components, including component coordinates, type,
  state, and properties. \\[0.4em]

  \code{/storage} &
  Per-user TeX files containing persistent ship data such as the ship
  name. \\

  \bottomrule
\end{tabular}
\end{center}

Even the database access was performed through TeX.

The custom SQL package wrote a query into a temporary file and then used
shell escape to invoke the MariaDB command-line client:

\begin{lstlisting}
\immediate\write18{
  mariadb ... < sqltmp.tex > sqloutput.tex
}
\end{lstlisting}

The result was then read back into the running TeX process.

Redis was handled in a similarly unusual fashion. TeX generated Redis
protocol messages and exchanged them with the Redis container through
a small helper script.

None of this was directly exploitable simply because it was unusual.

It did, however, mean that the boundaries between:

\begin{center}
  \displayfont
  HTTP
  $\longleftrightarrow$
  TeX
  $\longleftrightarrow$
  database
  $\longleftrightarrow$
  filesystem
  $\longleftrightarrow$
  shell
\end{center}

were extremely important.


% ------------------------------------------------------------
% 2.3
% ------------------------------------------------------------

\subsection{Persistent Ship Data}

Registration created more than a Redis account and a MariaDB ship entry.

LAMP also opened a file beneath \code{/storage} using the username and
stored the chosen ship name inside it:

\begin{lstlisting}
\immediate\openout\shipfile=/storage/\postusername{}
\store{\shipname}{\postshipname{}}
\immediate\closeout\shipfile{}
\end{lstlisting}

When the user later viewed their ship, the corresponding file was loaded
again:

\begin{lstlisting}
\input{/storage/\username{}.tex}
\end{lstlisting}

This detail becomes important later.

For now, the important observation is simply that attacker-controlled
data could survive a request, be written to disk as TeX, and then be
interpreted again during a future request.


% ------------------------------------------------------------
% 2.4
% ------------------------------------------------------------

\subsection{The First Useful Boundary}

The service contained a number of unusual trust boundaries, but the first
one that produced a useful exploitation primitive appeared very early in
request handling.

LAMP validated the path from the HTTP request line.

It then parsed the request headers.

One particular header could change the path that the application would
actually use afterwards.

That gave us our first primitive:

\begin{center}
  \displayfont\large\bfseries
  arbitrary file access through the \code{Path} header.
\end{center}